%! Author = stanislaw
%! Date = 22.03.2026

% Preamble
\documentclass[11pt]{article}

% Packages
\usepackage{amsmath}
\usepackage{tabularx}
\usepackage{booktabs}
\usepackage{seqsplit}
\usepackage{hyperref}

\title{Graph Isomorphisms and Automorphisms Algorithm Benchmarks}
\author{
    S. Piechota
    \and
    W. Walczak
    \and
    A. Zhuravlev
}
\date{\today}

% Document
\begin{document}

\maketitle
\tableofcontents
\newpage

\section{Color refinement}\label{sec:colorref}

    \subsection{Basic color refinement}\label{subsec:basic-color-refinement}
    \textbf{Basic color refinement} was introduced after hand-in of individual assignment in \textit{v1.0.0} of the project.

    \begin{table}[h!]
        \centering
        \small
        \begin{tabularx}{\textwidth}{|l|X|}
            \hline
            File name & \seqsplit{colorref 1.0.0} \\
            \hline
            test\_iter.grl & 8.44 \\
            colorref\_largeexample\_4\_1026.grl & 409.31 \\
            test\_empty.grl & 0.26 \\
            cref9vert3comp\_10\_27.grl & 10.11 \\
            test\_cycles.grl & 1.47 \\
            colorref\_smallexample\_4\_16.grl & 2.47 \\
            colorref\_smallexample\_2\_49.grl & 4.35 \\
            test\_cref9.grl & 2.05 \\
            test\_trees.grl & 2.29 \\
            colorref\_smallexample\_6\_15.grl & 5.23 \\
            colorref\_smallexample\_4\_7.grl & 0.94 \\
            lecture.grl & 1.39 \\
            test\_3reg.grl & 1.29 \\
            colorref\_largeexample\_6\_960.grl & 7897.13 \\
            \hline
        \end{tabularx}
        \caption{Color refinement algorithm runtimes (in milliseconds)}\label{tab:1}
    \end{table}

    \subsection{Fast color refinement}\label{subsec:fast-color-refinement}
    As clearly visible runtime for file \textit{colorref\_largeexample\_6\_960.grl} was clearly unsatisfactory (\hyperref[tab:1]{\textit{table 1}}).
    That is the reason that \textit{v1.1.0} introduced \textbf{fast color refinement}.

    \begin{table}[h!]
        \centering
        \small
        \begin{tabularx}{\textwidth}{|l|X|X|}
            \hline
            File name & \seqsplit{basic colorref 1.0.0} & \seqsplit{fast colorref 1.1.0} \\
            \hline
            colorref\_largeexample\_4\_1026.grl & 425.15 & 12866.37 \\
            cref9vert3comp\_10\_27.grl & 13.95 & 4.06 \\
            colorref\_largeexample\_6\_960.grl & 8407.62 & 13985.68 \\
            threepaths5.gr & 0.83 & 0.26 \\
            threepaths10.gr & 1.88 & 0.87 \\
            threepaths20.gr & 5.87 & 3.35 \\
            threepaths40.gr & 21.52 & 13.34 \\
            threepaths80.gr & 79.64 & 50.5 \\
            threepaths160.gr & 287.49 & 208.77 \\
            threepaths320.gr & 1135.27 & 868.5 \\
            threepaths640.gr & 4509.19 & 3752.91 \\
            \hline
        \end{tabularx}
        \caption{Basic color refinement and fast color refeinement runtime comparison (in milliseconds)}\label{tab:2}
    \end{table}

    In most cases fast color refinement gave better runtime that basic color refinement and in scenarios adapted for
    fast refinement gives better performance, but in large scenarios or edge
    cases (like \textit{colorref\_largeexample\_6\_960.grl}) fast color refinement could perform significantly worse than
    basic version (\hyperref[tab:2]{\textit{table 2}}, some entries from threepaths*.grl are excluded due to time limit).
    
    \subsection{Refactored color refinement}\label{subsec:refactored-fast-color-refinement}
    Version \textit{1.2.0} introduced more refactoring in the fast color refinement, which further improved the runtime
    (\hyperref[tab:3]{\textit{table 3}}, some entries from threepaths*.grl are excluded due to time limit).

    \begin{table}[h!]
        \centering
        \small
        \begin{tabularx}{\textwidth}{|l|X|X|}
            \hline
            File name & \seqsplit{fast colorref 1.1.0} & \seqsplit{fast colorref 1.2.0} \\
            \hline
            colorref\_largeexample\_4\_1026.grl & 12.42 & 0.26 \\
            colorref\_largeexample\_6\_960.grl & 13.81 & 0.12 \\
            threepaths5.gr & 0.0 & 0.0 \\
            threepaths10.gr & 0.0 & 0.0 \\
            threepaths20.gr & 0.0 & 0.0 \\
            threepaths40.gr & 0.01 & 0.0 \\
            threepaths80.gr & 0.05 & 0.01 \\
            threepaths160.gr & 0.21 & 0.02 \\
            threepaths320.gr & 0.88 & 0.02 \\
            threepaths640.gr & 3.77 & 0.05 \\
            \hline
        \end{tabularx}
        \caption{Fast refinement runtime in v1.1.0 and v1.2.0 (in seconds)}\label{tab:3}
    \end{table}

    As clearly visible, refactorization (based on changing the data type to set) brought extreme optimization in runtime
    that allowed to obtain satisfactory runtimes.

    Despite the great improvement as shown in \hyperref[tab:2]{\textit{table 2}} fast color refinement is susceptible to
    edge cases, that is why in version \textit{1.2.1} we made an attempt to refactorize basic color refinement as well
    (\hyperref[tab:4]{\textit{table 4}}).

    \begin{table}[h!]
        \small
        \begin{tabularx}{\textwidth}{|l|X|X|}
            \hline
            File name & \seqsplit{colorref 1.2.0} & \seqsplit{colorref 1.2.1} \\
            \hline
            test\_iter.grl & 3.91 & 1.37 \\
            colorref\_largeexample\_4\_1026.grl & 299.23 & 70.84 \\
            cref9vert3comp\_10\_27.grl & 4.89 & 1.66 \\
            colorref\_smallexample\_4\_16.grl & 1.24 & 0.44 \\
            colorref\_smallexample\_2\_49.grl & 2.72 & 0.81 \\
            test\_cref9.grl & 1.1 & 0.33 \\
            test\_trees.grl & 1.36 & 0.36 \\
            colorref\_smallexample\_6\_15.grl & 2.12 & 0.64 \\
            colorref\_largeexample\_6\_960.grl & 129.23 & 1689.66 \\
            threepaths5.gr & 0.49 & 0.13 \\
            threepaths10.gr & 0.8 & 0.27 \\
            threepaths20.gr & 1.62 & 0.87 \\
            threepaths40.gr & 3.6 & 2.9 \\
            threepaths80.gr & 14.7 & 11.07 \\
            threepaths160.gr & 12.79 & 43.94 \\
            threepaths320.gr & 25.74 & 175.45 \\
            threepaths640.gr & 54.79 & 716.55 \\
            threepaths1280.gr & 107.66 & 2890.96 \\
            threepaths2560.gr & 246.87 & 11949.6 \\
            \hline
        \end{tabularx}
        \caption{Color refinement runtime comparison v1.2.0 - v1.2.1 (in milliseconds)}\label{tab:4}
    \end{table}

    As predicted there exists singular scenarios where refactored basic color refinement was slightly better.
    But in general, with growing complexity, fast color refinement v1.2.0 was superior to all other attempts (\hyperref[tab:5]{\textit{table 5}}).

    \begin{table}[h!]
        \small
        \begin{tabularx}{\textwidth}{|l|X|}
            \hline
            File name & \seqsplit{fast colorref 1.2.0} \\
            \hline
            threepaths10240.gr & 2269.12 \\
            threepaths5120.gr & 692.83 \\
            threepaths2560.gr & 246.48 \\
            threepaths1280.gr & 109.68 \\
            threepaths640.gr & 53.24 \\
            threepaths320.gr & 25.53 \\
            threepaths160.gr & 13.06 \\
            threepaths80.gr & 6.54 \\
            threepaths40.gr & 3.37 \\
            threepaths20.gr & 1.68 \\
            threepaths10.gr & 0.92 \\
            threepaths5.gr & 0.55 \\
            wheelstar16.grl & 5.71 \\
            regulartwins.grl & 4.21 \\
            cubes9.grl & 32.38 \\
            wheeljoin25.grl & 2.17 \\
            wheeljoin19.grl & 2.81 \\
            cubes6.grl & 3.07 \\
            torus72.grl & 5.05 \\
            cographs1.grl & 1.91 \\
            torus144.grl & 15.38 \\
            wheelstar15.grl & 69.02 \\
            trees90.grl & 9.95 \\
            modulesD.grl & 45.84 \\
            products216.grl & 27.19 \\
            bigtrees3.grl & 14.05 \\
            \hline
        \end{tabularx}
        \caption{Color refinement v1.2.0 runtimes (in milliseconds)}\label{tab:5}
    \end{table}

\end{document}